\begin{section}{Ordered Sets}

An \textbf{ordered set} is a set $S$ together with a relation $<$ such that for any $x, y, z \in S$,
\begin{enumerate}[label=\textbf{(\Alph*)}]
    \item $x < y$ and $y < z$ implies $x < z$.
    \item Only one of the relations $x < y$, $x = y$, or $x > y$ holds.
\end{enumerate}
$x > y$ means $y < x$, $x \leq y$ means ($x < y$ or $x = y$) and $x \geq y$ means $y \leq x$.

\bigskip

\begin{remark}
    In an ordered set $S$, the relation $<$ is called a \textbf{order} (or a \textbf{total order}) on $S$
    and the set $S$ is called \textbf{totally ordered}. If a relation $<$ satisfies \textbf{(A)} only,
    then the relation $<$ is called a \textbf{partial order} on $S$ and the set $S$ is called \textbf{partially ordered}.
\end{remark}

\bigskip

\begin{example}
    The set $\N$ and $\Z$ are ordered sets, where $m > n \iff m - n \in \N$.
    The set $\Q$ is also an ordered set, where $m > n \iff m - n \in \Q^{+}$.
\end{example}

\bigskip

\begin{definition} \label{def:Boundedness}
    Let $S$ be an ordered set and $E \subseteq S$.
    \begin{itemize}
        \item $E$ is \textbf{bounded above} if there exists $M \in S$ such that $x \leq M$ for all $x \in E$.
        Such an $M$ is called an \textbf{upper bound} of $E$.
        \item $E$ is \textbf{bounded below} if there exists $m \in S$ such that $m \leq x$ for all $x \in E$.
        Such an $m$ is called a \textbf{lower bound} of $E$.
    \end{itemize}
\end{definition}

\bigskip

\begin{theorem} \label{thm:BoundedSubsetOfZHasExtremum}
    Any nonempty subset of $\Z$ which is bounded below has a least element,
    and any nonempty subset of $\Z$ which is bounded above has a greatest element.
\end{theorem}

\begin{proof}
    Let $E \subseteq \Z$ be nonempty and bounded below. Then there exists $m \in \Z$ such that $m \leq x$ for all $x \in E$.
    Let $F \coloneqq \setbuilder{x - m + 1}{x \in E} \subseteq \N$. By \Cref{thm:WellOrderingProperty},
    $F$ has a least element $n$.
    \[
    n \leq x - m + 1 \text{ for all } x \in E \implies n + m - 1 \leq x \text{ for all } x \in E.
    \]
    Since $n \in F$, there exists $x_{0} \in E$ such that $n = x_{0} - m + 1$,
    so $x_{0} = n + m - 1$ is the least element of $E$.
    \\[6mm]
    Let $S \subseteq \Z$ be nonempty and bounded above. Then there exists $M \in \Z$ such that $x \leq M$ for all $x \in S$.
    Let $T \coloneqq \setbuilder{-x + M + 1}{x \in S} \subseteq \N$. By \Cref{thm:WellOrderingProperty},
    $T$ has a least element $n$.
    \[
    n \leq -x + M + 1 \text{ for all } x \in S \implies M - n + 1 \geq x \text{ for all } x \in S.
    \]
    Since $n \in T$, there exists $x_{0} \in S$ such that $n = -x_{0} + M + 1$,
    so $x_{0} = M - n + 1$ is the greatest element of $S$.
\end{proof}

\bigskip

\begin{definition} \label{def:SupremumInfimum}
    Let $S$ be an ordered set and $E \subseteq S$. Suppose $E$ is bounded above.
    Suppose there exists an $\alpha \in S$ such that
    \begin{itemize}
        \item $\alpha$ is an upper bound of $E$.
        \item If $\beta < \alpha$, then $\beta$ is not an upper bound of $E$.
    \end{itemize}
    Then $\alpha$ is called the \emph{least upper bound} or \textbf{supremum} of $E$,
    denoted by
    \[
    \alpha = \sup E.
    \]
    Similarly, suppose $E$ is bounded below. Suppose there exists a $\alpha \in S$ such that
    \begin{itemize}
        \item $\alpha$ is a lower bound of $E$.
        \item If $\beta > \alpha$, then $\beta$ is not a lower bound of $E$.
    \end{itemize}
    Then $\alpha$ is called the \emph{greatest lower bound} or \textbf{infimum} of $E$,
    denoted by
    \[
    \alpha = \inf E.
    \]
\end{definition}

\bigskip

\begin{example}
    Let $S \coloneqq \setbuilder{p \in \Q^{+}}{p^{2} < 2}$. For any $p \in S$, let $q = p - \dfrac{p^{2} - 2}{p + 2}
    = \dfrac{2p + 2}{p + 2} > p$. Then $q^{2} - 2 = \dfrac{2(p^{2} - 2)}{(p + 2)^{2}} < 0$, so $q \in S$.
    This implies that any $p \in S$ is not an upper bound of $S$.
    \\[6mm]
    For any $p \in \Q^{+}$ with $p^{2} > 2$, $p$ is an upper bound of $S$. Let $q = p - \dfrac{p^{2} - 2}{p + 2}$ =
    $\dfrac{2p + 2}{p + 2} < p$. Then $q \in \Q^{+}$ and $q^{2} - 2 = \dfrac{2(p^{2} - 2)}{(p + 2)^{2}} > 0$, so $q$
    is an upper bound of $S$. This implies that such $p$ is not the least upper bound of $S$.
    \\[6mm]
    Since there exists no $p \in \Q$ such that $p^{2} = 2$, the set $S$ has no least upper bound in $\Q^{+}$.
\end{example}

\bigskip

\begin{definition} \label{def:LeastUpperBoundProperty}
    An ordered set $S$ is said to have the \textbf{least upper bound property}
    if every nonempty subset $E \subseteq S$ that is bounded above has a least upper bound in $S$.
\end{definition}

\bigskip

\begin{theorem} \label{thm:GreatestLowerBoundProperty}
    If $S$ is an ordered set with the least upper bound property, then any nonempty subset $E \subseteq S$
    that is bounded below has a greatest lower bound in $S$.
\end{theorem}

\begin{proof}
    Let $E$ be a nonempty subset of $S$ that is bounded below. Let $L$ be the set of all lower bounds of $E$.
    Then $L$ is nonempty and bounded above by any element of $E$. Let $\alpha = \sup L$. We claim that $\alpha = \inf E$.
    \\[6mm]
    Suppose $\alpha \notin L$. Then there exists $\beta \in E$ such that $\beta < \alpha$.
    Since $\alpha = \sup L$, there exists $\gamma \in L$ such that $\beta < \gamma \leq \alpha$.
    But $\gamma$ is a lower bound of $E$, so $\gamma \leq \beta$, a contradiction. Therefore, $\alpha \in L$.
    \\[6mm]
    Since $\alpha = \sup L$, for any $\delta > \alpha$, $\delta \notin L$. Therefore, $\alpha = \inf E$. 
\end{proof}

\end{section}